% GENERATED, edit the French unit or the translation table
\mxcorrections{1}\label{C01-CORR}
\textbf{Teacher use.} The answers below concern the tasks not worked in the lesson. First have students state their method; give graded help before the solution. The durations are estimates to be calibrated in class.
\subsection*{Diagnostic and remediation}
\begin{enumerate}
\item $11$: multiplication comes before addition. \textit{Action:} Rewrite the steps of two calculations with brackets.
\item $1$, $8$, $\dfrac{1}{4}$. \textit{Action:} Build the chain of powers by dividing by 2.
\item $80\times\num{1,05}=84$. \textit{Action:} Link \pc{5} to $\dfrac{5}{100}$ and then to the coefficient \num{1,05}.
\item $-2$ and $10$. \textit{Action:} Substitute x in brackets before calculating.
\item $10-3+1=8$. \textit{Action:} Write out the list, then recover the calculation.
\item $0,1,2,3$. \textit{Action:} Draw a two-column table, index and value; test range(1,4).
\end{enumerate}
\subsection*{Investigation activity}
Offer A: last week 45 kg, sum $\dfrac{6(20+45)}{2}=195$ kg. B: $\dfrac{20(1-\num{1,2}^6)}{1-\num{1,2}}=\num{198,5984}$ kg. B exceeds A in total, even though some of its early terms are smaller. A calculator is allowed here; argue from the cumulative total and not from the last term alone.
\subsection*{Immediate applications}
\begin{enumerate}
\item $2,3,6,11$. The formula $w_n=n^2+2$ works: it gives the first term and $(n+1)^2+2-(n^2+2)=2n+1$.
\item $37-17=5r$, so $r=4$ and then $u_0=17-3\times4=5$.
\item $u_{n+1}-u_n=-4\times3^n<0$: decreasing. $v_n$ alternates between 2 and $-2$: neither increasing nor decreasing.
\item $4+8+16+32=60$, or $\dfrac{4(1-2^4)}{1-2}=60$.
\item $u_3=6$: index 3 for $\leq6$, index 4 for $<6$.
\end{enumerate}
\subsection*{Exercises: method and results}
\paragraph{Exercise 1.1 -- Getting to grips with indices.} $u_0=0,u_1=-1,u_4=20$. Expand with brackets: $u_{n+1}=2(n+1)^2-3(n+1)=2n^2+n-1$; difference $4n-1$. The sign changes between indices 0 and 1: the sequence first decreases, then increases from index 1 onwards.
\paragraph{Exercise 1.2 -- Choosing the right model.} $a_n=100+15n$ and $b_n=100(\num{1,1})^n$, so $a_2=130$, $b_2=121$. A assumes a constant net addition; B a constant net rate. In reality, rounding and capacity limit these models.
\paragraph{Exercise 1.3 -- Recovering a sequence.} $u_3=q^2u_1$ so $q^2=9$. Positivity forces $q=3$, hence $u_0=2$ and $u_5=486$. Without this condition, $q=-3$ would also be possible, with $u_0=-2$ and the same $u_5=486$.
\paragraph{Exercise 1.4 -- Rows of seats.} The last row has $18+11\times2=40$ seats. Sum $\dfrac{12(18+40)}{2}=348$. Check: average 29 seats per row, between 18 and 40.
\paragraph{Exercise 1.5 -- Accumulating squares.} We obtain 4 then 9 squares. A square of integer side n completed by a row and a column needs $2n+1$ cells. The formula $n^2$ has the right first term and $ (n+1)^2-n^2=2n+1$. The generation is the same.
\paragraph{Exercise 1.6 -- Two repayments.} For A: $a_n=600-50n$ for $0\leq n\leq12$, then 0. For B: $b_n=\dfrac{600}{2^n}$. $b_5=\num{18,75}$ and $b_6=\num{9,375}$; strict threshold at index 6. B never gives exactly 0 in this ideal model: do not announce a finite repayment.
\paragraph{Exercise 1.7 -- Reading a list.} The first three elements are 2, 5, 14. An initial value and five additions give six elements; indices 0 to 5. The last one is $u_5=365$.
\paragraph{Exercise 1.8 -- Delayed review.} $u_7=12+5(-3)=-3$. Six terms: $\dfrac{6(12-3)}{2}=27$. The 5 counts the passes, the 6 the terms.
\paragraph{CQFD, exercise 9 p.~64, questions a and b.} a: $-1,0,3,8$. b: $2,\dfrac72,4,\dfrac{17}4$. The first index is 0; the denominator of b is $n+1$ and never vanishes on $\N$.
\paragraph{CQFD, exercise 11 p.~64, all questions.} Starting from $u_0=2$, the rule $u_{n+1}=2u_n-1$ gives $u_1=3$, $u_2=5$, $u_3=9$, $u_4=17$. The rule uses the previous term, not its index.
\subsection*{MCQ: answers and reasons}
\textbf{MCQ 1.}
\begin{enumerate}
\item A. $3\times4+2=14$.
\item C. Three additions: $2+3\times5=17$.
\item D. The coefficient is $1+\dfrac{8}{100}$.
\item B. $9-3+1=7$.
\item C. $\dfrac{20\times 21}{2}=210$.
\end{enumerate}
\textbf{MCQ 2.}
\begin{enumerate}
\item D. The difference is $-3\times2^n<0$.
\item A. Four terms; $\dfrac{1-2^4}{1-2}=15$.
\item C. $u_3=2$, $u_4=1$.
\item B. The terms alternate between two distinct values.
\item D. Eleven indices, from 0 to 10.
\end{enumerate}
\subsection*{True or false}
\begin{enumerate}
\item False: $u_n=-2^n$ is strictly decreasing.
\item False: $u_n=0$ and $v_n=n(n-1)(n-2)$ coincide at indices 0, 1, 2, but not at index 3.
\item True: each term equals 1 and there are n+1 terms.
\end{enumerate}
\paragraph{Differentiated teaching.} For a technical block, offer only the first step (domain, factor or reference sample space), then withdraw the help. For an independent student, ask for a second method and the treatment of a limiting case. Orally, allow an initial reformulation in the language the student has mastered, then require a precise mathematical sentence in French.
